\section{Baboon Utopia}

\begin{quotex}
It is therefore evident that a nation affirming a hierarchical civilization that does not recognize the separation of the human from the divine, is naturally moral, even ignoring the same terms of morality: its existence is the law of itself, hence it consciously creates its own life, gives soul and form to its nation, making it a ``model'' nation.
\flright{\textsc{Massimo Scaligero}, \emph{La Razza di Roma}}
\end{quotex}


\paragraph{Preface}


Tradition accepts the findings of science insofar as they are restricted to their own domain, but not when science attempts to go beyond its legitimate methods. Here are some examples.

\textbf{Man is under many laws}. Man is subject to the laws of physics and chemistry, as well as biology. Moreover, man is subject to the laws of genetics including the law of the survival of the fittest. Behaviour that does promote reproductive success will lead to the extinction of the group, tribe, or nation.

\textbf{Higher laws}. The laws mentioned above are recognized by science. However, there are higher laws, such as the psychological impulses, cultural norms, sociological and political laws. Of course, one strives to transcend these softer laws, but ultimately the harder laws of science.

\textbf{Man has an animal soul}. I often read about either the intelligence of animals or how animal-like human beings are. There should be no surprise in this, since Tradition recognizes that the human being also has an animal soul (as well as a vegetable soul).

\paragraph{Introduction}


I recently discovered the popular biologist \textbf{Robert Sapolsky}, who consistently gets several hundred thousand views on youtube. To his credit, he is a competent and engaging speaker, even if he is an atheist and positivist. Apparently, the educated elite who watch him need to be reassured that they are nothing but naked primates, unlike the rubes who believe that there is something special about humans.

Joe Rogan interviews Sapolsky\footnote{\url{https://www.youtube.com/watch?v=obmt_PkIfBE}} about chimpanzees, baboons, and even human beings. After dealing with Rogan's morbid interest in parasites, we get to the relevant part about the baboons.

\paragraph{Baboon Life}


Baboons live in small troops that are hierarchical and patriarchal. Such a social organisation has served the baboons well for their two million year history. Nevertheless, this concerns Sapolsky as he claims it leads to stress in the baboons. Sapolsky also seems to harbour an antipathy to hierarchy and, as we shall see, approves of its opposite.

On the other hand, Tradition approves of hierarchy and sees that the negation of hierarchy leads to conflict. Apparently, the baboons resist the natural hierarchy in the troop, causing fights and disruption.

Another characteristic is that adolescent baboons leave their home troop and bond with another. This leads to a circulation of genetic material.

\paragraph{The Forest Troop}


One troop, named the Forest Troop, demonstrated a behavioural change due to a unique event. The more aggressive males stopped hunting when they discovered a garbage dump from a nearby resort. Unfortunately, some of the meat was tainted and the baboons dies from tuberculosis. That left the Forest Troop with the less aggressive males and a surfeit of females.

Sapolsky was intrigued by the behavioural changes within the troop. The males become epicene, since they groomed each other more often, and the females more promiscuous by approaching new males within days. Sapolsky calls this a ``cultural change'' since it doesn't depend on genetic changes, but, rather, it is learned behaviour. Specifically, newly arrived adolescents adopted the new behaviour simply by repeating what they saw. I don't know why this was such a surprise, since we all know about monkey see monkey do\footnote{\url{https://www.urbandictionary.com/define.php?term=monkey\%20see\%20monkey\%20do}}.

Of course, Sapolsky is thrilled about baboon utopia and believe it shows hope that humans can make a similar cultural change from the patriarchy. Maybe you don't look forward to picking bugs off your guy friends, but the upside is that the women will be friendlier to you. Obviously, once the females are no longer dominated by their monkey, they are free to exhibit their sexuality more openly. Jan and Dean even sang about such a utopia in \textit{Two girls for every boy}\footnote{\url{https://www.youtube.com/watch?v=5cvtD2U5mvM}}.

So Sapolsky's hope for human cultural change is effeminate males and aggressive females. In my opinion, there was no cultural change, but rather the same genetic programming in an altered social condition. Even if true, a cultural change is still subject to the law of the survival of the fittest. The Grand Experiment of the Forest Troop did not last very long, and their behaviour reverted back to patriarchy and hierarchy. That may be the real lesson for human revolutionaries.

\paragraph{Free will vs Power}


Sapolsky then tackles the notion of Free Will, which he denies. Unfortunately, he conflates free will and power. He believes that the inability to control every one of our movements refutes free will, when it simply means that we don't have the power to control. The oddball example he gives is that of an epileptic who hits someone during an episode. He claims, without proof, that in the Middle Ages, the epileptic would have been prosecuted for assault. Even at that time, the notion of intention was understood, even if Sapolsky does not. The epileptic was unable to control his movements, so there was no crime; certainly, that is not an example that refutes free will.

At the very least, we have some control over our thoughts. It is absurd to believe that our genes, or hormones, are sophisticated enough to decide which TV program to watch tonight.

\paragraph{A Reductionist Explanation}


In the end, Sapolsky cannot avoid moral judgments, despite his denial of free will. Obviously, he despises the hierarchy of the Middle Ages. However, were we so inclined, we could likewise provide a purely biological and genetic explanation of the Medieval leaders. Specifically, they served to:

\begin{itemize}


\item Enforce group unity

\item Promote order

\item Cull society of genetic defects, e.g., psychopaths
\end{itemize}


The Medievals promoted group unity through their common worship. As history has shown, religious and ideological conflicts can greatly disrupt human societies. By eliminating such conflicts before they could become widespread, they promoted genetic success.

Unlike the baboons that were in conflict over access to food and to females, the Medievals enforced monogamous mating. This gave nearly all the males access to females, greatly reducing male rivalries.

Through the common use of capital punishment, the unanticipated effect was to remove psychopaths and incorrigible criminals from the gene pool. Moreover, capital punishment served a higher purpose, as described in \textit{Medieval Executions: The View from the Scaffold}\footnote{\url{https://www.medievalists.net/2017/11/medieval-executions-view-scaffold/}}.


\flrightit{Posted on 2020-02-11 by Cologero}
